% A game that shows what you can write in clues and responses.
\documentclass{article}
% Which PDF to make (the web game ignores this):
%   questions  clues only, e.g. a study sheet for students
%   answers    clues and responses
%   solutions  clues, responses, and worked solutions
\usepackage[solutions]{../jeopardy}

\gametitle{Formatting Demo}

% Macros defined here work inside math in every clue.
\newcommand{\R}{\mathbb{R}}
\newcommand{\dd}[1]{\frac{d}{d#1}}
\newcommand{\abs}[1]{\left|#1\right|}
\DeclareMathOperator{\Span}{span}

\begin{document}

\begin{category}{Inline math}
  \begin{clue}
    Inline math with dollar signs: $f(x) = x^3 - 2x$.
  \response
    $f'(x) = 3x^2 - 2$
  \end{clue}
  \begin{clue}
    Inline math with parentheses: \(\sin^2\theta + \cos^2\theta\).
  \response
    \(1\)
  \end{clue}
  \begin{clue}
    A macro from the preamble: $\dd{x}\left(e^{x}\right)$
  \response
    $e^x$
  \end{clue}
  \begin{clue}
    Another macro: $\abs{-7}$ and the reals $\R$.
  \response
    $7$
  \end{clue}
  \begin{clue}
    A declared operator: $\Span\{v_1, v_2\}$
  \response
    All linear combinations $a v_1 + b v_2$
  \end{clue}
\end{category}

\begin{category}{Display math}
  \begin{clue}
    Display math on its own line:
    \[ \int_0^{\pi} \sin x \, dx \]
  \response
    $2$
  \solution
    An antiderivative of $\sin x$ is $-\cos x$, so
    \begin{align*}
      \int_0^{\pi} \sin x \, dx &= \Big[-\cos x\Big]_0^{\pi} \\
        &= -\cos\pi + \cos 0 \\
        &= 1 + 1 = 2.
    \end{align*}
  \end{clue}
  \begin{clue}
    Double-dollar display:
    $$ \lim_{n \to \infty} \left(1 + \frac{1}{n}\right)^n $$
  \response
    $e$
  \end{clue}
  \begin{clue}
    An \texttt{align*} environment:
    \begin{align*}
      x + y &= 5 \\
      x - y &= 1
    \end{align*}
  \response
    $x = 3$, $y = 2$
  \solution
    Add the equations: $2x = 6$, so $x = 3$. Then $y = 5 - x = 2$.
  \end{clue}
  \begin{clue}[dailydouble]
    A matrix: $A = \begin{pmatrix} 1 & 2 \\ 3 & 4 \end{pmatrix}$. Find $\det A$.
  \response
    $-2$
  \solution
    $\det A = (1)(4) - (2)(3) = 4 - 6 = -2$
  \end{clue}
  \begin{clue}
    A piecewise function:
    \[ f(x) = \begin{cases} x^2 & x \ge 0 \\ -x & x < 0 \end{cases} \]
  \response
    $f(-3) = 3$
  \end{clue}
\end{category}

\begin{category}{Text formatting}
  \begin{clue}
    \textbf{Bold}, \emph{emphasis}, and \textit{italics}.
  \response
    Formatting works.
  \end{clue}
  \begin{clue}
    A dollar sign in text: it costs \$5.
  \response
    \$5
  \end{clue}
  \begin{clue}
    Line breaks:\\ second line\\ third line
  \response
    Three lines
  \end{clue}
  \begin{clue}
    Quotes and dashes: ``quoted'' text, pages 3--5, and a pause---like this.
  \response
    Typographic quotes and dashes
  \end{clue}
  \begin{clue}
    Accents: Poincar\'e, Erd\H{o}s, G\"odel.
  \response
    Names with accents
  \end{clue}
\end{category}

\begin{category}{Lists}
  \begin{clue}
    Which of these is prime?
    \begin{itemize}
      \item $21$
      \item $23$
      \item $25$
    \end{itemize}
  \response
    $23$
  \end{clue}
  \begin{clue}
    Put in order:
    \begin{enumerate}
      \item Differentiate
      \item Set equal to zero
      \item Solve
    \end{enumerate}
  \response
    Finding critical points
  \end{clue}
  \begin{clue}
    Short clue.
  \response
    Short response.
  \end{clue}
  \begin{clue}
    A longer clue to show how the board shrinks text to fit. Let
    $f:\R \to \R$ be continuous on $[a,b]$ and differentiable on $(a,b)$.
    Then there is some $c \in (a,b)$ where the tangent line is parallel
    to the secant line through the endpoints, that is,
    \[ f'(c) = \frac{f(b) - f(a)}{b - a}. \]
  \response
    The Mean Value Theorem
  \end{clue}
  \begin{clue}
    Comments are ignored. % This comment won't appear.
  \response
    Nothing after the percent sign shows.
  \end{clue}
\end{category}

\begin{category}{Math in a name $\Sigma$}
  \begin{clue}
    Category names can contain math.
  \response
    Response
  \end{clue}
  \begin{clue}
    Greek: $\alpha, \beta, \gamma, \delta, \varepsilon$
  \response
    Response
  \end{clue}
  \begin{clue}
    Sets: $A \cup B$, $A \cap B$, $A \setminus B$
  \response
    Response
  \end{clue}
  \begin{clue}
    Logic: $\forall x\, \exists y : y > x$
  \response
    Response
  \end{clue}
  \begin{clue}
    Blackboard bold: $\mathbb{N} \subset \mathbb{Z} \subset \mathbb{Q} \subset \R$
  \response
    Response
  \end{clue}
\end{category}

\begin{final}{Famous Constants}
  It equals $\displaystyle\sum_{n=0}^{\infty} \frac{1}{n!}$.
\response
  What is $e$?
\solution
  The Taylor series of $e^x$ is $\displaystyle\sum_{n=0}^{\infty} \frac{x^n}{n!}$.
  Setting $x = 1$ gives $e$.
\end{final}

\end{document}
